\documentclass[11pt]{article}
\usepackage[T1]{fontenc}
\usepackage{lmodern}
\usepackage[margin=1in]{geometry}
\usepackage{amsmath,amssymb,amsthm,mathtools}
\usepackage{microtype,booktabs,tabularx,array,enumitem,fancyvrb}
\usepackage[dvipsnames]{xcolor}
\usepackage[colorlinks=true,linkcolor=MidnightBlue,citecolor=MidnightBlue,urlcolor=MidnightBlue]{hyperref}
\hypersetup{pdftitle={A scalar parity refinement and a conditional faster SAT contraction},pdfauthor={Research note prepared with ChatGPT for Vladimir Reshetnikov}}
\setlength{\emergencystretch}{3em}
\setlength{\parskip}{0.35em}
\setlength{\parindent}{0pt}
\setlist{itemsep=0.25em,topsep=0.4em}
\allowdisplaybreaks[2]
\newtheorem{theorem}{Theorem}[section]
\newtheorem{lemma}[theorem]{Lemma}
\newtheorem{proposition}[theorem]{Proposition}
\newtheorem{corollary}[theorem]{Corollary}
\theoremstyle{definition}
\newtheorem{definition}[theorem]{Definition}
\newtheorem{assumption}[theorem]{Unproved assumption}
\theoremstyle{remark}
\newtheorem{remark}[theorem]{Remark}
\newcommand{\poly}{\operatorname{poly}}
\newcommand{\supp}{\operatorname{supp}}
\newcommand{\F}{\mathbb F}
\newcommand{\Z}{\mathbb Z}
\newcommand{\N}{\mathbb N}
\newcommand{\ind}{\mathbf 1}
\newcommand{\ceil}[1]{\left\lceil #1\right\rceil}
\newcommand{\floor}[1]{\left\lfloor #1\right\rfloor}
\newcommand{\upstream}{7ef5c0cc5ca4b0e072b7a050fca09ab24bc6fa3a}
\newcommand{\sourcebase}{https://github.com/openai/math/blob/7ef5c0cc5ca4b0e072b7a050fca09ab24bc6fa3a/preprints/The-Exponential-Time-Hypothesis-is-False-October-25-2026/build/}
\title{A Scalar Parity Refinement of the\newline Subexponential SAT Counter\\[0.5em]
\large Proved arithmetic improvements and a conditional route\newline to a better running-time exponent}
\author{Research note prepared for Vladimir Reshetnikov}
\date{October 26, 2026}
\begin{document}
\maketitle
\begin{abstract}
We examine the October 25 revision of OpenAI's preprint \emph{The
Exponential Time Hypothesis is False}, pinned to the repository revision
specified below. We obtain an elementary, exact replacement for its dense-query
parity lift. The replacement removes the truncated polynomial variable $V$,
uses scalar residues, and improves the retained incidence-degree cap from
$k(2n+a)$ to $k(n+a)$ for $a$ sparse degree-$k$ tests. Its correctness,
coefficient divisibility, and sharpness are proved and checked by executable
finite tests. This gives precise improvements in arithmetic representation and
local operations, but preserves the manuscript's leading general-input exponent
$O_k(n\log\log n/\sqrt{\log n})$.

For a stronger improvement we propose applying contraction a second time to
a compact encoding of its mode sum. We derive the conditional bound
$(n+m)^{O_k(1)}2^{O_k(n(\log\log n)^{3/2}/\log n)}$, including the
necessary precision analysis and retuning of two preprocessing stages. The
missing step is an explicit compact factorization of the resulting high-degree
group tensors. We state that step as an unproved assumption and explain why
neither the existing theorem nor index compression supplies it. Thus this note
does not claim a newly proved improved leading worst-case SAT exponent.
\end{abstract}

\begin{quote}
\textbf{Status of the requested breakthrough.} A concrete arithmetic refinement
is proved. A stronger asymptotic algorithm is proposed conditionally. The
structural factorization needed to turn that proposal into a theorem has not
been constructed. The included Python programs validate the finite algebra and
indexing identities; they are not complete subexponential SAT solvers.
\end{quote}

\newpage
{\small\setlength{\parskip}{0pt}\tableofcontents}
\newpage

\section{Result, baseline, and scope}

\subsection{The comparison that matters}
Let $n$ be the number of declared variables, $m$ the number of input clauses,
and $k$ a fixed clause width. The user's case is $k=3$. Put
$\lambda=\ceil{\log_2 n}$ and $\nu=\ceil{\log_2\lambda}$ when $n$ is
sufficiently large. All asymptotic statements allow an explicit exhaustive
fallback on a bounded range of smaller sizes. Logarithm bases affect constants
only.

The manuscript states a deterministic exact counter with polynomial space and
the first two bounds in Table~\ref{tab:bounds}. These are stronger than the
bound described on the repository's Lean scope page. The latter is
$2^{O_k(n(\log\log n)^2/\sqrt{\log n})}$ times a polynomial in encoded
input length~\cite{source,scope}. This distinction matters: removing one
$\log\log n$ from that weaker description would not improve the revised
manuscript.

\begin{table}[ht]
\centering\small
\begin{tabularx}{\textwidth}{@{}p{0.25\textwidth}p{0.39\textwidth}X@{}}
\toprule
Method or setting & Exponent, omitting polynomial factors & Status here \\
\midrule
Source, arbitrary density & $O_k(n\log\log n/\sqrt{\log n})$ & Inherited manuscript theorem \\
Source, $m\le Dn$ & $O_{k,D}(n/\sqrt{\log n})$ & Inherited manuscript theorem \\
Scalar parity replacement & $O_k(n\log\log n/\sqrt{\log n})$ & Proved replacement; same leading order \\
Second contraction, arbitrary density & $O_k(n(\log\log n)^{3/2}/\log n)$ & Conditional on Assumption~\ref{ass:LF} \\
Second contraction, $m\le Dn$ & $O_{k,D}(n\sqrt{\log\log n}/\log n)$ & Conditional on Assumption~\ref{ass:LF} \\
Union-support compiler & $O_k(n/\sqrt{\log n})$ & Separate unproved route \\
\bottomrule
\end{tabularx}
\caption{Every algorithm in the table would be deterministic and use
polynomial space under its stated assumptions. A conditional row is not an
established improvement.}\label{tab:bounds}
\end{table}

The proposed second-contraction exponent is strictly smaller in asymptotic
order: the ratio to the manuscript's arbitrary-density exponent is
\[
 \frac{n(\log\log n)^{3/2}/\log n}
      {n\log\log n/\sqrt{\log n}}
 =\sqrt{\frac{\log\log n}{\log n}}\longrightarrow0.
\]
This comparison concerns the bounding functions. It is neither an observed
speedup nor a proof that the unproved factorization exists.

\subsection{Source provenance and validation boundaries}
The inspected repository commit is
\begin{center}\texttt{\upstream}.\end{center}
The principal files are the main manuscript, its quantitative reduction, and
its quantitative contraction proof~\cite{source,reduction,contraction}. A
machine-readable manifest in the package records the exact retrieved URLs and
hashes. We read the algorithms and their parameter arguments, rather than
assuming the title alone establishes their correctness.

The elementary algebra in Sections~\ref{sec:scalar}--\ref{sec:resources} is
proved independently below. Claims that it gives a complete subexponential
counter invoke the source contraction and reduction theorems. We did not
rebuild the complete Lean formalization, implement the full contraction
engine, or independently verify every theorem in the preprint. The repository's
comparator statement is a specification; it should not be confused with an
executed proof check. No claim of priority is made for the familiar parity
indicator identity itself.

\section{Locating the two asymptotic losses}\label{sec:bottlenecks}

The quantitative reduction first obtains systems of $a=O_k(n)$ polynomial
equations over $\F_2$. Each equation has bounded degree and
$(\log n)^{O_k(1)}$ listed monomials. It gives every test--variable incidence
its own Boolean indeterminate. If $N$ is the total number of these incidences,
then $N\le n(\log n)^{O_k(1)}$.

Its polynomial lift has surviving total incidence degree $K=O_k(n)$. After
grouping tests and variables, it enumerates degree allocations and encodes
selected incidence subsets by binary indices. This is the construction labeled
\texttt{eq:sharp-wire-budget} in the source~\cite{reduction}. Its size estimate is
\begin{equation}\label{eq:wire}
 M\le J^2+\log_2\binom{N+K}{K}
 \le J^2+K\log_2\!\frac{\mathrm e(N+K)}{K}
 =O_k(n\log\log n),\qquad J=\ceil{n^{1/4}}.
\end{equation}
Our first refinement changes $K$ and the coefficient ring, but leaves the order
of this bound unchanged.

The second loss is inside the contraction engine. With $M$ active wires and
$u$ retained valuation levels, the source has at most $O(M/D)$ modes of each
positive weight and a working precision $P=O(MD/u)$. Counting weighted mode
tuples gives
\begin{equation}\label{eq:partition-source}
 \exp\!\left(O\!\sqrt{\frac MD\,P}\right)
 =\exp(O(M/\sqrt u)).
\end{equation}
The available $u$ is of order $\log s$ at the root of an $s$-wire contraction.
Its recursive costs sum to $\exp(O(s/\sqrt{\log s}))$. Combining this with
\eqref{eq:wire} produces the dense-input exponent. The candidate improvements
in this note address these two steps explicitly.

\section{An exact scalar parity replacement}\label{sec:scalar}

\subsection{Normalization and the local indicator}
Write a normalized polynomial over $\F_2$ as
\[
 P_i(x)=c_i+\sum_{j=1}^{t_i}m_{ij}(x),\qquad c_i\in\{0,1\},
\]
where each $m_{ij}$ is a distinct nonconstant squarefree monomial of degree
at most $k$. Repeated monomials have been canceled modulo two. Replace
$x_v$ in test $i$ by the incidence variable $z_{iv}$, and write $z_i$ for
that test's variables. Define
\begin{equation}\label{eq:E}
 E_i(z_i)=\frac{1+(-1)^{c_i}
                    \prod_{j=1}^{t_i}(1-2m_{ij}(z_i))}{2}.
\end{equation}
The division is performed over the integers before reducing modulo a power of
two. It is not inversion of $2$ in $\Z/2^p\Z$.

\begin{lemma}[Exact indicator]\label{lem:indicator}
The expression \eqref{eq:E} is an integer polynomial. For every Boolean
assignment $y$ to its incidence variables,
\[
 E_i(y)=\ind_{\{P_i(y)=0\}}.
\]
\end{lemma}
\begin{proof}
The constant term of the numerator is $1+(-1)^{c_i}$, which is even.
Every nonempty selection of $r$ factors from the product has coefficient
$(-1)^{c_i+r}2^r$, also even. On Boolean inputs, each $m_{ij}(y)$ is
zero or one, so $1-2m_{ij}(y)=(-1)^{m_{ij}(y)}$. Thus the product with
the sign $(-1)^{c_i}$ equals $(-1)^{P_i(y)}$, and the claimed indicator
follows.
\end{proof}

Let $\overline E_i$ denote the multilinear representative obtained by
replacing each positive power of an incidence variable by its first power.
It still agrees with $E_i$ on the Boolean cube. For $S\subseteq I_i$, use
$z_{i,S}=\prod_{v\in S}z_{iv}$ and write
$c_{i,S}=[z_{i,S}]\overline E_i$.

For example, if $P=x_1x_2+x_2x_3$, then
\[
 \overline E=1-x_1x_2-x_2x_3+2x_1x_2x_3.
\]
The cubic coefficient already records a factor of two. If instead the two
monomials are disjoint cubic monomials $A,B$, the indicator is
$1-A-B+2AB$. The coefficient at incidence degree six has valuation one.
These are the elementary phenomena used in the general degree filtration.

\subsection{Coefficientwise divisibility}
We use $v_2(0)=+\infty$.

\begin{lemma}[Local degree versus valuation]\label{lem:valuation}
For every nonempty $S\subseteq I_i$,
\begin{equation}\label{eq:local-val}
 v_2(c_{i,S})\ge
 \max\left\{0,\ceil{\frac{|S|}{k}}-1\right\}.
\end{equation}
The coefficient for $S=\varnothing$ is integral as well.
\end{lemma}
\begin{proof}
After dividing by two, a term selecting $r\ge1$ monomials in
\eqref{eq:E} has coefficient $\pm2^{r-1}$. Multilinearizing that term
gives the union of the selected supports, of size at most $kr$. A contribution
to the coefficient indexed by $S$ therefore has
$r\ge\ceil{|S|/k}$. Each contribution has the stated divisibility, so their
sum does too. The statement for the empty support follows from
Lemma~\ref{lem:indicator}.
\end{proof}

Work now in $R_p=\Z/2^p\Z$, and form
\[
 F(z)=\prod_{i=1}^{a}\overline E_i(z_i)\quad\text{in }R_p[z].
\]
Different tests have disjoint incidence variables, so this product is
multilinear. If a monomial uses support $S_i$ in each test, its coefficient is
exactly $\prod_i c_{i,S_i}$. No expansion over all supports is required to
evaluate that particular coefficient.

\begin{theorem}[Sharper incidence-degree cap]\label{thm:cap}
An incidence monomial of total degree $d$ in $F$ has coefficient divisible by
\[
 2^{\max\{0,\ceil{d/k}-a\}}.
\]
Consequently, modulo $2^p$, every incidence monomial of degree greater than
\begin{equation}\label{eq:cap}
 K'=k(p+a-1)
\end{equation}
vanishes. For exact counting on $n$ declared variables, take $p=n+1$, giving
$K'=k(n+a)$.
\end{theorem}
\begin{proof}
Let $b\le a$ be the number of nonempty supports $S_i$. Summing
\eqref{eq:local-val} over those supports gives at least
\[
 \sum_{S_i\ne\varnothing}\ceil{|S_i|/k}-b
 \ge\ceil{d/k}-a.
\]
Valuations are also nonnegative. If $d>k(p+a-1)$, the last expression is
at least $p$, so the coefficient is zero in $R_p$.
\end{proof}

The source bound is $k(a+2(p-1))$, whereas the new bound is
$k(a+p-1)$~\cite{reduction}. For the at-most-$6n$-test systems of the
general reduction, these upper caps are respectively $8kn$ and $7kn$.
For $k=3$ this particular worst-case comparison is $24n$ versus $21n$.
The exact improvement depends on the actual number of tests; these numbers
are bounds, not the number of variables of a new SAT instance.

\begin{proposition}[Sharpness in this representation]\label{prop:sharp}
The cap $k(p+a-1)$ cannot be decreased for arbitrary such test systems merely
by strengthening a universal coefficientwise degree argument.
\end{proposition}
\begin{proof}
In test $i$, take $r_i\ge1$ mutually disjoint monomials, all of degree $k$,
and take no constant term. On independent incidence variables the coefficient
selecting every one of these monomials in every test is
\[
 (-1)^{\sum_i r_i}2^{\sum_i r_i-a}.
\]
Choose $\sum_i r_i=p+a-1$. Its incidence degree is $k(p+a-1)$ and its
valuation is $p-1$, so it survives modulo $2^p$. Tests may reuse original
variables; only their incidence copies are independent. The assertion concerns
the representation, not a lower bound for solving the equations.
\end{proof}

\section{Integrating the replacement into the exact counter}\label{sec:integration}

\subsection{The counting identity}
Identify the incidence copies by setting $z_{iv}=x_v$. By
Lemma~\ref{lem:indicator}, over the integers,
\begin{equation}\label{eq:exact-count}
 N(\mathcal P)=\sum_{x\in\{0,1\}^n}\prod_{i=1}^aE_i(x)
             =\sum_xF((z_{iv}=x_v)_{iv})\pmod{2^{n+1}}.
\end{equation}
The count lies between zero and $2^n$, so its least nonnegative residue modulo
$2^{n+1}$ recovers the exact answer. Variables absent from every test remain
declared; their free choices must be counted. Empty test lists and constant
equations are handled before the main construction.

Use the source's test grouping, variable grouping, incidence-degree arrays,
and subset indices, with the cap $K'$ from Theorem~\ref{thm:cap}. For a
specified support $S_i$ in each test, a test-group tensor has scalar entry
\begin{equation}\label{eq:test-entry}
 \prod_{i\in\mathcal T_g} c_{i,S_i}\pmod{2^{n+1}}.
\end{equation}
For a variable group, count one factor of two for each variable whose incidence
copies are all absent from the selected monomial. All other variables are
forced to one when summing that monomial. Invalid binary subset indices receive
entry zero. These are the same combinatorial counting rules as in the source;
the coefficients are now scalar.

\subsection{Coefficient evaluation without expanding the lift}
The formal product in \eqref{eq:E} can have exponentially many monomials.
The intended entry routine does not construct that product.

\begin{lemma}[Scalar coefficient oracle]\label{lem:oracle}
For every $S\subseteq I_i$,
\begin{equation}\label{eq:mobius}
 c_{i,S}=\sum_{T\subseteq S}(-1)^{|S|-|T|}
                       \ind_{\{P_i(\mathbf1_T)=0\}},
\end{equation}
where variables outside $T$ are set to zero. A coefficient can be computed in
$2^{|S|}\poly(t_i,k,|I_i|,p)$ time and polynomial space, without $V$ or
the source's polynomials $e_j$.
\end{lemma}
\begin{proof}
For a multilinear polynomial $G$, evaluation at $\mathbf1_T$ is the sum
of its coefficients indexed by subsets of $T$. Subset-lattice inversion gives
\eqref{eq:mobius}, and Lemma~\ref{lem:indicator} supplies its values. Evaluate
the sparse polynomial over $\F_2$ for each $T$, add or subtract one, and
retain one accumulator. No table of all subsets is needed. Its absolute value
is at most $2^{|S|}$ before reduction modulo $2^p$.
\end{proof}

The sparse-query inputs have $|I_i|\le k\lambda^c$ for a fixed $c$.
Thus the coefficient oracle takes $2^{\lambda^{O_{k,c}(1)}}$ time, within
the permitted $2^{O(s^{9/10})}$ local-entry bound when $s=n\nu$.
The high-incidence branch and grouping bound the maximum degree by
$n^{3/4}\lambda^{O_{k,c}(1)}\le s^{9/10}$ for sufficiently large $n$.
All tensor descriptions remain polynomial in $n$.

\begin{corollary}[Valid replacement in the source algorithm]\label{cor:replacement}
Assuming the source's contraction theorem and preprocessing lemmas, replacing
its dense-query lift by \eqref{eq:E} gives a deterministic exact counter with
polynomial space and the same proved leading arbitrary-density bound
\[
 (n+m)^{O_k(1)}2^{O_k(n\log\log n/\sqrt{\log n})}.
\]
Its dense-query contractions can use $\ell=1$ and scalar initial coefficients.
\end{corollary}
\begin{proof}
The preceding counting identity and tensor construction establish correctness.
Equation~\eqref{eq:wire} still holds with $K'=\Theta_{k,c}(n)$. Every
other hypothesis of the source theorem is unchanged or weakened, and it
explicitly permits $\ell=1$. The source's exact sparsification and greedy
parity-test selection then use this valid query oracle.
\end{proof}

\section{Precisely what becomes faster or smaller}\label{sec:resources}

Let $p=n+1$. The original dense-query coefficient ring is
$(\Z/2^p\Z)[V]/(V^p)$. The replacement starts in $\Z/2^p\Z$.
The contraction proof still uses its unramified extension of degree $e$ and
working precision $q$; neither of those is removed.

\begin{table}[ht]
\centering\small
\begin{tabularx}{\textwidth}{@{}p{0.38\textwidth}XX@{}}
\toprule
Quantity & Original lift & Scalar lift \\
\midrule
Root coefficient payload & $p^2$ bits & $p$ bits \\
One coefficient at extension degree $e$, precision $q$ & $eqp$ bits & $eq$ bits \\
Coefficient products in one schoolbook $V$-multiplication & $p(p+1)/2$ & $1$ \\
Retained incidence-degree bound & $k(a+2p-2)$ & $k(a+p-1)$ \\
Local subset-evaluation output & A vector of $p$ coefficients & One Boolean parity value \\
Readout after contractions & Coefficient-sum map after truncated products & Scalar residue \\
\bottomrule
\end{tabularx}
\caption{Exact representation and subroutine comparisons. These are not
implementation-independent speedup factors for the entire solver.}
\end{table}

The coefficient payload improves by a factor of $p$, and schoolbook
multiplication in the removed axis improves by $\Theta(p^2)$. Fast polynomial
multiplication would give a different arithmetic comparison. The whole program
also stores graphs, mode tuples, interpolation data, and immutable entry
descriptions, so a factor-$n$ reduction in total space is not asserted.
Likewise, no factor-$n^2$ reduction in total runtime follows just from the
schoolbook multiplication count.

This replacement is \emph{not} justified by substituting $V=1$ into the old
quotient ring. That substitution is not a ring homomorphism: $V^p=0$ there,
whereas $1^p=1$. The new construction uses an independent exact integer
indicator and proves coefficientwise divisibility before reducing modulo
$2^p$. This removes the source's need to control a coefficient-sum map that
does not preserve products.

An implementation should use \eqref{eq:mobius} for entry lookups, retain
the existing creation-precision and canonical-lift rules, and avoid storing
the full expanded indicator. It may return zero immediately if
$\ceil{|S|/k}-1\ge p$. It can multiply the per-test scalar coefficients
sequentially. These changes are concrete even before any stronger
factorization theorem is available.

\section{Why several immediate exponent improvements fail}\label{sec:barriers}

\subsection{Incidence-support indexing has a real entropy cost}
Take $a$ tests, each with $t$ disjoint nonconstant degree-$k$ monomials, and no
constant. On independent incidence variables,
\[
 \overline E_i\equiv1+P_i\pmod2.
\]
Their product has exactly $(t+1)^a$ distinct odd-coefficient incidence
monomials. Each test independently contributes its constant or one of its $t$
monomials; distinct incidence choices cannot coincide. Every one of these
terms survives at every positive precision.

For $a=\Theta(n)$ and $t=(\log n)^c$, an explicit injective index of all
these supports needs $\Omega(n\log\log n)$ bits. There are only
$2^{O(\sqrt n\log n)}$ degree arrays in the source construction. By the
pigeonhole principle, distributing the supports among those arrays leaves an
array requiring $\Omega(n\log\log n)$ index bits. The tighter cap in
Theorem~\ref{thm:cap} does not remove these odd terms.

This is a limitation of a representation, not of SAT. The same example with
linear tests can be solved by Gaussian elimination. A successful improvement
must use additional algebraic structure, combine incidence selections, or
change the contraction method.

\subsection{The weighted-partition estimate is sharp for its relaxation}
Consider an abstract set with $A$ modes of each weight $1,2,\ldots$ and a
budget $P$, in the regime $P/A\to\infty$. Let $q=\floor{c\sqrt{P/A}}$ for a sufficiently small positive
constant $c$. Choosing each of the $Aq$ modes of weights at most $q$ with
exponent zero or one produces at least $2^{Aq}$ tuples. Even the tuple choosing
all of them consumes only
\[
 A\frac{q(q+1)}2<P.
\]
Hence this relaxation contains $\exp(\Omega(\sqrt{AP}))$ tuples.
Under the source scaling, their ordinary degree is only
$Aq=\Theta(M/\sqrt u)\le M$. Imposing that ordinary-degree bound does not
improve the exponent. Keeping an additional starting valuation of order
$\log s$ is likewise insufficient when the important weights are of order
$D/\sqrt u\gg\log s$.

This does not prove that every such tuple has a nonzero coefficient in every
actual network. It shows why a better estimate of the currently enumerated
relaxation alone is not a new algorithm.

\subsection{Nonnegative Hermite tapers have a restricted limitation}
Here is a more explicit reason to look beyond taper tuning. Let nonnegative
multiplicities $q_i$ vanish for $i\ge u$, and suppose their degree coverage is
$\sum_{i<u}2^{-i}q_i\ge cM$. Define
\[
 L_i=2q_i-\sum_{j\ge1}\frac{j}{2^j}q_{i+j}.
\]
If $K=\max_iL_i$, then
\[
 q_i\le K/2+\mathbb E[q_{i+J}],\qquad
 \Pr(J=j)=\frac{j}{2^{j+1}}.
\]
Because $J\ge1$, iterating reaches the zero boundary in at most $u-i$ steps.
Thus $q_i\le K(u-i)/2$. Also
$\max_iq_i\ge cM/2$, so $K\ge cM/u$.

For the finite grid of the source with $a\ge2u$, truncating the sum at
$j=a-i$ changes the loss by at most
$O(u2^{-u})\max_iq_i$. Combining this with
$K\ge2\max_iq_i/u$ preserves the $\Omega(M/u)$ lower scale for large
$u$. This concerns the cardinal Hermite precision analysis with nonnegative
multiplicities. It does not rule out a different recovery method that exploits
dependencies between jets.

\subsection{Other changes that do not certify a faster bound}
Using the source clause network immediately after sparsification is weaker:
for $k=3$ its cost has exponent proportional to
$n[\sigma+B/\sqrt{\log n}]$, where
$B=O((\sigma^{-1}\log(2/\sigma))^2)$. Ignoring secondary logarithms, its
best balance is $n/(\log n)^{1/6}$, worse than the source's dense bound.

Changing the counting modulus does not remove the dominant Hermite loss
$O(MD/u)$. At an odd prime larger than the interpolation grid, denominators
become units, but the binary geometric tile ratio no longer gains a power of
that prime with each mode order. The existing truncation proof therefore
cannot simply be reused. Replacing powers by a Newton basis introduces
$v_2(b!)=b-O(\log b)$, which changes linear valuation growth only by a
constant factor.

\section{The stronger proposal: contract the mode sum again}\label{sec:conditional}

The source already expresses a deformed contraction as
\begin{equation}\label{eq:mode-sum}
 F(T)=\sum_{\beta:\,\sum_Zw_Z\beta_Z<P}
                         \prod_{g=1}^r Z_{g,\beta}\pmod{2^P},
\end{equation}
where every mode has one source group and one target group. Instead of
enumerating every $\beta$ separately, encode bundles of such modes as edges
of a new network. The index compression itself is valid. The missing
factorization appears only after it.

\subsection{One deformation stage with more groups}
Start with a source-theorem input at ambient size $s$, precision $O(s)$,
ambient degree bound $\Delta=\ceil{s^{9/10}}$ on the input maximum
degree, and local entries evaluable in
$\poly(s)\exp(O(s^{9/10}))$ time and polynomial space. For clarity, take
polynomial-size descriptions, as in the SAT applications. Put
\begin{equation}\label{eq:first-stage}
 \lambda=\ceil{\log_2s},\quad r=\ceil{\lambda},\quad
 u=\floor{\frac{\lambda}{4096\log_2\lambda}},\quad
 a=2u,\quad L=2^a,\quad D_{\rm arr}=r^{2a}.
\end{equation}
These parameters are used for one stage, not for a recursively growing-group
algorithm. If $M\le s/\lambda$, evaluate directly. Otherwise $M>s/\lambda$.

Writing $c=1/4096$, we have
\[
 D_{\rm arr}=s^{4c+o(1)},\qquad L=s^{o(1)},\qquad
 2D_{\rm arr}^{2}\Delta=s^{9/10+8c+o(1)}=o(s/\lambda).
\]
The source's other packing inequalities $M\ge LD_{\rm arr}$ and
$2L\le\Delta$ also hold for large $s$. Its polynomial-time labeling and
packet construction therefore remain available~\cite{contraction}. The
child contraction for group $g$ has at most $6M/r$ edges.

At an active interpolation node of valuation $i<u$, put $D=L/2^i$. The
tapered precision and mode multiplicity estimates remain
\begin{equation}\label{eq:one-stage}
 P=O(MD/u),\qquad A_{\rm total}\le6M/D,
 \qquad P\le s^{1+o(1)}.
\end{equation}
Indeed $D\ge2^{u+1}$ absorbs both the inherited $O(s)$ precision and the
$O(Ma)$ recovery term. All parameters and all intermediate precisions are
polynomial in $s$.

For fixed $\beta$, evaluate $Z_{g,\beta}$ by direct summation of its
$O(s/r)$ internal bits. The original entries and the new packet-entry
coefficients give the bound
\begin{equation}\label{eq:group-time}
 E(s)=\poly(s)\exp(O(s/r+s^{9/10})),
\end{equation}
with polynomial space. The packet-entry term
$\exp(O(L\log s))$ is absorbed. This argument does not assume that a
fast child contraction automatically works at the inflated precision.

\subsection{Exact pair budgets and a shorter index network}
Bundle modes by their unordered pair $e$ of distinct source and target groups;
retain their directions in the descriptions. Let $A_e$ be a bound on the
number of its modes at each weight, chosen so that
$\sum_eA_e\le6M/D$. Enumerate \emph{exact} budgets $b_e\ge0$ with
$\sum_eb_e<P$. There are at most
\begin{equation}\label{eq:budget-arrays}
 (P+1)^{r(r-1)/2}=\exp(O((\log s)^3))
\end{equation}
arrays. Exact budgets are essential: using an upper bound for each pair and
then summing over all bounds would count many tuples repeatedly.

Let $R_e(b_e)$ count the tuples in bundle $e$ of exact weight $b_e$.
For positive budgets and nonempty bundles, the weighted-partition estimate
implies $\log R_e(b_e)=O(\sqrt{A_eb_e})$. A zero budget has one tuple.
Skip an impossible budget. Index a feasible fiber using
$\ceil{\log_2R_e(b_e)}$ bits, assigning zero to unused binary words.
Then the number of index wires satisfies
\begin{align}
 W_{\rm array}
 &\le r^2+O\!\left(\sum_e\sqrt{A_eb_e}\right)\notag\\
 &\le r^2+O\!\sqrt{\left(\sum_eA_e\right)\left(\sum_eb_e\right)}
 =O(M/\sqrt u+r^2)\notag\\
 &=O\!\left(s\sqrt{\frac{\log\log s}{\log s}}\right).
 \label{eq:compressed-wires}
\end{align}
Use a sufficiently large fixed constant to choose a common ambient upper
budget
\[
 W=C_0\ceil{s\sqrt{\log\log s/\log s}}.
\]
Here $W$ is an upper bound; some arrays have far fewer wires. Counting,
ranking, and unranking weighted tuples are polynomial in $s$, because the
numeric budget $P$, number of modes, and integer bit lengths are polynomial
in $s$. Appendix~\ref{app:rank} states a concrete dynamic program.

For each array, the new graph has one tensor per group, computing
$Z_{g,\beta}$ from the decoded incident indices. Its contraction is exactly
the part of \eqref{eq:mode-sum} with that array. Summing arrays restores the
whole expression. This exact reformulation, including invalid-index zeros,
is the part tested by the second prototype.

\subsection{The missing structural result}
\begin{assumption}[Compact factorization of the group tensors, LF]\label{ass:LF}
For the family of input networks under consideration, there is a uniform
algorithm which, for every feasible first-stage interpolation instance and
exact budget array, factors each group tensor as a Boolean tensor network
such that:
\begin{enumerate}[label=(\roman*)]
\item Each group tensor is preserved for every assignment to its boundary
index bits, including zeros at invalid binary indices. The equality holds in
the current coefficient algebra modulo $2^P$.
\item The combined factorizations have $O(W)$ total edges, including all
new internal edges, and maximum degree at most $\ceil{W^{9/10}}$.
They have polynomially many vertices and polynomial-size descriptions.
\item Construction and every local factor-entry evaluation take at most
$\poly(s)\exp(O(s/r+s^{9/10}))$ time and polynomial space.
\item Entries are fixed at their creation precision $P$ and admit the same
canonical lifts used in the source. The original unramified coefficient field
is retained.
\end{enumerate}
A constant-factor subdivision of edges to make the result bipartite is allowed.
\end{assumption}

\textbf{Assumption LF is not proved in this note or supplied by the source.}
For the SAT consequence, it suffices to prove it for the networks generated
by the clause and sparse-query reductions. Assuming it for all source-theorem
networks is a stronger research target.

The obstruction is concrete. The unfactored graph has only $r\asymp\log s$
group tensors. In the regime $W_{\rm array}\asymp W$, some degree is at
least $2W_{\rm array}/r\gg W^{9/10}$. A short program for an entry is not
a factorization with small degrees. Restoring the original child networks
retains about $2M$ original row and column wires in total, which also fails
the required $O(W)$ edge budget.

Applying the existing packing argument directly to the high-degree graph
does not give another $\sqrt{\log W}$ saving. If the second group alphabet
has fixed size $R$, its packing condition is
$W\ge2R^{4a'}\Delta'$. Since $\Delta'\ge cW/r$, it permits only
$a'=O(\log r)$ useful levels. Meanwhile the first stage has
$u=O(\log s/\log r)$. Their product is only $O(\log s)$; the apparent
two-stage saving returns to the original square-root-logarithm scale.
LF must exploit extra structure of the group tensors, beyond index length.

\section{Conditional running time and all remaining costs}\label{sec:bound}

\subsection{A modest extension of the existing contraction analysis}
The second contraction has initial precision $P=s^{1+o(1)}$, which may exceed
its wire budget $W$. This requires an explicit check; the original initial
precision hypothesis cannot be quoted unchanged.

\begin{lemma}[Precision and entry-oracle extension]\label{lem:extension}
Let $S$ be a size parameter and retain the source's constant
$\epsilon=1/(128\cdot1024)$. Its contraction proof extends to initial
precision
\[
 P_0\le S^{1+\epsilon/2},
\]
in a fixed unramified coefficient ring of degree $O(\log S)$ with at least
$S^4$ residue elements and polynomial-time arithmetic in its given integral
basis, as in the source construction. A polynomial truncation length is permitted. If an
original tensor entry takes $\poly(S)\exp(B(S))$ time and polynomial
space, with fixed creation-precision semantics, the total time is
\[
 \poly(S)\exp\bigl(O(S/\sqrt{\log S})+B(S)\bigr),
\]
with polynomial space, under the other source size and degree hypotheses.
\end{lemma}
\begin{proof}
Let $\lambda=\ceil{\log_2S}$ and use the source's parameters
$u_d=\floor{\epsilon\lambda/2^d}$ and
$a_d=2\sum_{j\le d}u_j$. At an active root,
$M>S/\sqrt\lambda$, and every retained node has
$D\ge2^{u_0+1}\ge cS^\epsilon$. Therefore
\[
 \frac{P_0}{MD/u_0}
 \le O(\lambda^{3/2}S^{-\epsilon/2})=o(1).
\]
The needed bound $P_t=O(MD/u_0)$ still holds. At depths $d\ge1$,
$a_{d-1}\ge2u_0$, so $S2^{a_{d-1}}$ absorbs the extra initial
precision. The source's inherited-precision bound and subsequent argument
then apply. Its original first invariant at depth zero is thus replaced,
not silently assumed.

All precisions remain polynomial in $S$. Its branching and interpolation
counts are unchanged. The cost of an original entry is a multiplicative
factor, because descriptions are immutable and evaluated directly at their
creation precision. Newly generated packet entries have their existing
$\exp(O(S^{9/10}))$ bound, absorbed by the main exponent. The same
depth-first traversal gives polynomial space. The proof's interpolation uses
an adequate residue field, rather than requiring a freshly constructed
extension at each application. Polynomial truncation length changes only
polynomial arithmetic factors.
\end{proof}

In our application $P=W^{1+o(1)}$, so the lemma's explicit precision
inequality eventually holds. Reuse the first field, whose degree is
$4\ceil{\log_2s}=O(\log W)$ and whose size is at least $s^4>W^4$
for sufficiently large $s$. Tensoring two unrelated extensions and assuming
the result is a field would not be a valid shortcut.

\subsection{The conditional contraction theorem}
\begin{theorem}[Conditional second contraction]\label{thm:conditional}
For a family satisfying Assumption~\ref{ass:LF}, the source-theorem input
networks can be contracted deterministically in polynomial space and time
\begin{equation}\label{eq:improved-engine}
 \poly(s)\,2^{O(s\sqrt{\log\log s}/\log s)}.
\end{equation}
\end{theorem}
\begin{proof}
Perform the one-stage deformation of Section~\ref{sec:conditional}, then
apply LF to each exact-budget graph. By Lemma~\ref{lem:extension}, its second
contraction takes time
\[
 \poly(s)\exp\left(O\!\left(\frac W{\sqrt{\log W}}\right)
                          +O(s/r+s^{9/10})\right).
\]
Since $r=\Theta(\log s)$ and
$W=O(s\sqrt{\log\log s/\log s})$, this is
$\poly(s)\exp(O(s\sqrt{\log\log s}/\log s))$.
The factor \eqref{eq:budget-arrays} and all polynomially many interpolation
instances are smaller-order multiplicative overheads. Hermite recovery gives
the original contraction. Process arrays, interpolation values, and recursive
calls sequentially; every stage uses polynomial space. Explicit parameter
checks and exhaustive handling of bounded exceptional sizes give a uniform
conditional algorithm.
\end{proof}

\subsection{Retuning preprocessing is necessary}
Substituting $s=n\nu$ into \eqref{eq:improved-engine} suggests the target
$n\nu^{3/2}/\lambda$, but that substitution alone is incomplete. Two
branches of the source use $2^{O(n/\sqrt\lambda)}$ time. Either would
eventually dominate the proposed bound if left unchanged.

Use the sparsification parameter
\[
 \sigma=1/\ceil{\log_2n},
\]
in place of order $1/\sqrt\lambda$. The source's sparsification lemma
then gives $2^{O(n/\lambda)}$ branches and occurrence bound
\[
 B=O_k((\lambda\log\lambda)^{k-1}).
\]
The sparse-check row lengths remain fixed powers of $\lambda$: the source
has $P_{\rm hash}=20000B^2\mu^2$ and $D_0=320000B^2\mu^3$, with
$\mu=O(\log n)$ after normalization. Candidate enumeration is still
$2^{\lambda^{O_k(1)}}$, and every query still has at most $6n$ tests of
degree at most $k$ and polylogarithmic monomial count~\cite{source}.

In the query reduction, also replace the high-incidence cutoff by
\begin{equation}\label{eq:cutoff}
 D_* =\max\{1,\ceil{N_0/n}\}\ceil{\log_2n},
\end{equation}
where $N_0$ is the initial incidence count. At most $n/\lambda$ variables
exceed this cutoff. Fixing them costs $2^{O(n/\lambda)}$ branches. All
remaining degrees are still polylogarithmic, the test and variable groups
still have degree $n^{3/4}\lambda^{O_k(1)}$, and the wire budget is still
$O_k(n\nu)$. The degree arrays contribute
$2^{O_k(\sqrt n\log n)}$, also smaller than the target.

\begin{corollary}[Conditional improvement for fixed-width SAT and counting]
\label{cor:conditional-sat}
If LF holds for the networks arising from these reductions, exact counting
and hence $k$-SAT are deterministic and polynomial-space with time
\[
 (n+m)^{O_k(1)}
 2^{O_k(n(\log\log n)^{3/2}/\log n)}.
\]
For $m\le Dn$ with fixed $D$, the direct clause network gives
\[
 (n+m)^{O_{k,D}(1)}
 2^{O_{k,D}(n\sqrt{\log\log n}/\log n)}.
\]
\end{corollary}
\begin{proof}
The dense case uses $s=\Theta(n\nu)$ and the retuned preprocessing just
verified. Its additional exponents $n/\lambda$, $\sqrt n\log n$, and
$\lambda^{O_k(1)}$ are absorbed by $n\nu^{3/2}/\lambda$. At fixed
density, the source clause network has $s=\Theta_{k,D}(n)$ and needs no
dense reduction. Apply Theorem~\ref{thm:conditional}.
\end{proof}

These are conditional mathematical implications, not certified new running
times. The needed construction is LF. It would be a substantial structural
result, and no claim that it is true follows from its usefulness.

\section{A narrower alternative: aggregate equal variable unions}\label{sec:union}

There is a separate target for removing only the dense-case $\log\log n$.
In the scalar polynomial $F(z)$, identify all copies of variable $v$ and
reduce by $x_v^2=x_v$. For $S\subseteq[n]$ define
\[
 c_S=\sum_{\substack{Q\text{ an incidence support}\\
                         \operatorname{vars}(Q)=S}}[z_Q]F(z)
                      \pmod{2^{n+1}}.
\]
Then
\begin{equation}\label{eq:union}
 N(\mathcal P)\equiv\sum_{S\subseteq[n]}2^{n-|S|}c_S
                          \pmod{2^{n+1}}.
\end{equation}
This follows by summing each resulting squarefree monomial over the Boolean
cube. There are only $2^n$ possible unions, even when the incidence expansion
has $2^{\Theta(n\log\log n)}$ supports.

If the functional \eqref{eq:union} can be uniformly compiled into an
$O(n)$-wire network of maximum degree at most $n^{9/10}$, with polynomial
descriptions and $2^{O(n^{9/10})}$-time, polynomial-space entry routines,
the existing contraction theorem gives
\[
 (n+m)^{O_k(1)}2^{O_k(n/\sqrt{\log n})}.
\]
The source's original preprocessing costs fit that target. This proposal
does not require the improved engine, but it also lacks a compiler theorem.

The coefficient identity is not that theorem. An $n$-port tensor for $c_S$
has forbidden degree, and directly evaluating a coefficient by a subset sum
can require $2^{|S|}$ work. The prototype's dictionary of union supports
is an exponential reference calculation for small inputs. Its correctness
does not certify a fast compiler. A useful next investigation would seek a
factorization that preserves the quotient $x_v^2=x_v$ without explicitly
retaining all test--variable choices.

\section{Implementation artifacts and checks}\label{sec:checks}

\subsection{Scalar parity prototype}
\texttt{code/parity\_lift.py} represents a multilinear polynomial by a
dictionary from support bit masks to integer coefficients. Its reference
multiplication uses support union, reflecting $x_v^2=x_v$. It implements
normalization, the exact parity product, direct coefficient inversion, the
old truncated-$V$ comparison, and independent exhaustive common-zero counts.
Only the coefficient oracle is intended as a replacement primitive; expanding
whole dictionaries is used to validate small examples.

The executed suite, with seed $20261026$, passed:
\begin{center}\small
\begin{tabular}{@{}lr@{}}
\toprule
Check & Cases \\
\midrule
Exhaustive normalized polynomials & 2,304 \\
Truth-table comparisons & 34,816 \\
Coefficient and divisibility comparisons & 34,816 \\
Random systems compared with assignment enumeration & 300 \\
Old truncated-$V$ products versus scalar products & 250 \\
Global incidence-degree checks & 120 \\
Sharp degree-cap examples & 12 \\
\bottomrule
\end{tabular}
\end{center}
Constants, empty equations, unused variables, and cancellation of repeated
monomials are included. The old-versus-new test checks products before the
source readout, rather than assuming that summing coefficient vectors is a
ring homomorphism.

\subsection{Weighted-mode indexing prototype}
\texttt{code/weighted\_partitions.py} counts, ranks, and unranks nonnegative
exponent tuples of prescribed exact weight. It also checks a three-group mode
sum by comparing direct enumeration with exact pair-budget arrays and
binary indices. Invalid padding words have weight zero in the encoded sum.
The toy group functions contain nonlinear terms and depend only on incident
modes, so the test exercises the relevant factorization of the summation.

The executed suite passed 40 weight lists, 360 exact fibers, 1,525 exact
rank round trips, 400 cumulative budgets, 5,065 cumulative round trips,
360 invalid-rank checks, and 12 modular grouped-sum identities. At strict
budget six, the direct sum has 297 tuples; the encoded calculation has 297
valid rank words across 462 exact budget arrays and five invalid padding
words. Results agree modulo both 97 and 101. These small moduli are test
arithmetic only; the proposed contraction still uses the source's two-adic
coefficient ring.

\subsection{What was and was not validated}
The tests validate exact identities and the finite indexing construction.
They do not validate LF, the asymptotic contraction implementation, or a new
SAT runtime. The elementary proofs establish the general algebraic claims;
finite tests are supplementary. The source theorem is an external dependency,
and the reported formalization scope is attributed to the repository.

The archive contains the self-contained LaTeX source and compiled PDF, both
standard-library prototypes, recorded JSON results, a source manifest, a
machine-readable research-status file, and build/verification scripts. Run
\begin{Verbatim}[fontsize=\small]
python3 code/parity_lift.py
python3 code/weighted_partitions.py
python3 verify_package.py
python3 build_report.py
\end{Verbatim}
The final command requires a TeX installation. Timing fields in the JSON
describe the small validation suite and are not benchmarks of a SAT solver.

\section{Recommended next mathematical step}

The scalar replacement can be implemented and formalized independently of
the speculative compression work. Its central proof obligations are short:
the parity product identity, support-degree divisibility, coefficient
inversion, and the existing tensor-network counting identity. It also removes
the auxiliary coefficient-sum map from the dense-query proof interface.

For the requested leading exponent improvement, the most specific next task
is to construct or refute LF on the group tensors generated by sparse
degree-three parity queries. The construction must expose every auxiliary
wire, count their total, and provide bounded-degree local entries with fixed
creation-precision semantics. An equality of formal sums or a fast oracle for
one large tensor is insufficient. Small exact tensor-rank experiments may
help choose a representation, but do not replace these size and uniformity
bounds.

If LF fails on that family, the union-support compiler in
Section~\ref{sec:union} is a distinct route to a smaller improvement. Either
success would bypass a demonstrated representation cost. Until such a result
is proved, the strongest established conclusion of this report remains the
scalar arithmetic refinement, with the source's leading exponent unchanged.

\appendix
\section{Exact weighted-tuple indexing}\label{app:rank}

For positive weights $w_0,\ldots,w_{q-1}$ and an integer maximum total $B$,
let $D_i(b)$ count tuples on coordinates $i,\ldots,q-1$ of exact weight $b$.
Use boundary values $D_q(0)=1$ and $D_q(b)=0$ for $b>0$. Then
\[
 D_i(b)=D_{i+1}(b)+
 \begin{cases}
 D_i(b-w_i),&b\ge w_i,\\
 0,&b<w_i.
 \end{cases}
\]
Compute $i$ in decreasing order and $b$ in increasing order. This uses
$O(qB)$ integer cells and arithmetic operations; it is pseudopolynomial in
the numeric budget. In the application $q$ and $B=P$ are polynomial in $s$,
so it is polynomial in the input size parameter used in the proof.

For lexicographic ranking in an exact fiber, choosing the next exponent
$j$ leaves $D_{i+1}(b-jw_i)$ completions. Sum these counts for all smaller
choices. Unranking selects the unique interval containing the rank. Cumulative
strict-budget indexing first chooses the total $b<P$ using the counts
$D_0(b)$ and then uses the exact-fiber procedure. A binary word whose integer
value is at least the fiber size is invalid. This gives a bijection, rather
than an encoding that accidentally counts several upper budgets for one
tuple.

\section{Source locations for an implementation patch}\label{app:locations}

The source uses labels rather than stable theorem numbers across revisions.
These labels identify the relevant locations at the pinned commit:
\begin{center}\small
\begin{tabularx}{\textwidth}{@{}p{0.36\textwidth}X@{}}
\toprule
Source label or file & Relevance \\
\midrule
\texttt{lem:rational-idempotent},\newline\texttt{lem:total-order} & The old lift and auxiliary $V$-truncation; replaced only in the quantitative dense-query route. \\
\texttt{lem:tensor-sparse-query} & Install the scalar coefficient oracle and counting identity here. \\
\texttt{eq:sharp-incidence-degree} & Replace the cap by $k(n+a)$. \\
\texttt{eq:sharp-wire-budget} & Keep the subset-index proof with the new cap; its asymptotic order is unchanged. \\
\texttt{eq:tensor-precision-bound} & Preserve canonical lifts; extend the root analysis only for the conditional second contraction. \\
\texttt{eq:tensor-mode-multiplicity},\newline\texttt{eq:tensor-partition-count} & The weighted-mode index bounds used in the proposed compression. \\
\texttt{eq:tensor-recursion-identity} & The sum regrouped by exact pair budgets. \\
\texttt{lem:sparsify},\newline\texttt{lem:sparse-checks},\newline\texttt{lem:greedy-hashes} & Retune sparsification to $1/\lambda$ only for the conditional stronger bound. \\
\texttt{lean/docs/378.md} & The repository's stated formalization scope, weaker than the revised quantitative manuscript. \\
\bottomrule
\end{tabularx}
\end{center}

\begin{thebibliography}{9}
\bibitem{source}
OpenAI, \emph{The Exponential Time Hypothesis is False}, October 17, 2026,
revised October 25, 2026. Inspected at repository commit
\texttt{\upstream}.
\href{https://github.com/openai/math/blob/7ef5c0cc5ca4b0e072b7a050fca09ab24bc6fa3a/preprints/The-Exponential-Time-Hypothesis-is-False-October-25-2026/build/main_full_version.tex}{Pinned main LaTeX source}.

\bibitem{reduction}
OpenAI, \emph{A sharper bound for exact counting}, quantitative reduction
appendix of~\cite{source}.
\href{https://github.com/openai/math/blob/7ef5c0cc5ca4b0e072b7a050fca09ab24bc6fa3a/preprints/The-Exponential-Time-Hypothesis-is-False-October-25-2026/build/quantitative-reduction.tex}{Pinned quantitative-reduction.tex}.

\bibitem{contraction}
OpenAI, \emph{Proof of the contraction theorem}, quantitative contraction
appendix of~\cite{source}.
\href{https://github.com/openai/math/blob/7ef5c0cc5ca4b0e072b7a050fca09ab24bc6fa3a/preprints/The-Exponential-Time-Hypothesis-is-False-October-25-2026/build/quantitative-contraction.tex}{Pinned quantitative-contraction.tex}.

\bibitem{scope}
OpenAI, \emph{The Exponential Time Hypothesis is False: Lean formalization
scope}, repository documentation, inspected at the same commit.
\href{https://github.com/openai/math/blob/7ef5c0cc5ca4b0e072b7a050fca09ab24bc6fa3a/lean/docs/378.md}{Pinned scope statement}.
\end{thebibliography}
\end{document}
